\chapter*{Glossary}
\addcontentsline{toc}{chapter}{Glossary}
\markboth{Glossary}{Glossary}
\index{glossary}%

This glossary grows with each chapter and collects the cross-cloud data-engineering terms
introduced so far. Definitions are platform-neutral unless a platform is named explicitly.

\begin{description}
  \item[Access tier] A storage pricing class that trades retrieval cost and latency for lower
    at-rest price (for example Azure Hot, Cool, Cold, and Archive). See also \emph{storage class}.
  \item[ACID transaction] Atomicity, consistency, isolation, and durability guarantees for table
    writes; Delta Lake adds ACID semantics to Parquet files via its transaction log.
  \item[AlloyDB] Google Cloud's PostgreSQL-compatible engine with a columnar accelerator for
    analytical queries, positioned as the cloud-native tier above Cloud SQL.
  \item[Apache Spark] A distributed data-processing engine; Databricks is its original commercial
    steward, with counterparts in AWS Glue/EMR, Azure Synapse/Fabric, and GCP Dataproc.
  \item[Asset Bundle] Databricks' infrastructure-as-code format (\texttt{databricks.yml}) defining
    jobs, pipelines, and deployment targets for CI/CD.
  \item[Auto Loader] Databricks' incremental file-ingestion source (\texttt{cloudFiles}) with
    checkpointed exactly-once semantics, schema inference/evolution, and rescued-data handling.
  \item[Aurora] Amazon's cloud-native relational engine with distributed storage, Global Databases,
    and a serverless (v2) compute option.
  \item[Availability SLA] The contractual uptime a provider guarantees for a service, distinct from
    \emph{durability}, which concerns data loss rather than reachability.
  \item[Bigtable] Google Cloud's wide-column, high-throughput, low-latency NoSQL store for
    time-series, IoT, and ad-tech workloads; exposes the HBase API and scales by nodes and clusters.
  \item[Bronze / Silver / Gold] The medallion layers: raw append-only capture, validated and
    conformed tables, and curated dimensional or aggregate tables for consumption.
  \item[Bucket (storage)] The top-level container for objects in AWS S3 and Google Cloud Storage;
    Azure calls the equivalent a \emph{container}; Unity Catalog governs file access via volumes.
  \item[Change Data Capture (CDC)] Streaming inserts, updates, and deletes from source transaction
    logs to targets; on Databricks applied with \texttt{MERGE INTO} or DLT \texttt{AUTO CDC}.
  \item[Checkpoint] The state directory that makes Structured Streaming and Auto Loader exactly-once
    across restarts; also a compacted Parquet snapshot of the Delta transaction log.
  \item[Clustering key] A user-defined column set that Snowflake uses to co-locate related rows within
    micro-partitions, improving partition pruning on selective queries; compare Delta liquid clustering.
  \item[CMEK / CMK] Customer-managed encryption keys (GCP CMEK; Azure customer-managed keys via Key
    Vault), giving the customer control over the key that encrypts data at rest.
  \item[Cosmos DB] Azure's globally distributed multi-model NoSQL database, billed in Request Units
    (RU) and supporting multiple APIs (NoSQL, MongoDB, Cassandra).
  \item[DBU] Databricks Unit, the compute billing unit consumed per second of cluster or serverless
    runtime; total cost is DBU rate plus underlying cloud VM/infrastructure charges.
  \item[Delta Lake] The open table format (created by Databricks) adding ACID transactions, time
    travel, schema enforcement, and deletion vectors to Parquet files in cloud object storage.
  \item[Delta Live Tables (DLT)] Databricks' declarative pipeline framework with streaming tables,
    materialized views, expectations, and an event log; successor concept: Lakeflow Declarative Pipelines.
  \item[Delta Sharing] An open protocol for sharing live Delta tables to any recipient without
    copying data or granting storage credentials.
  \item[Deletion vector] A Delta Lake row-level delete marker that avoids rewriting entire Parquet
    files for \texttt{DELETE}, \texttt{UPDATE}, and \texttt{MERGE} operations.
  \item[DynamoDB] Amazon's managed key-value and document NoSQL database, with capacity measured in
    read/write capacity units (RCU/WCU) or on-demand, and Global Tables for multi-region.
  \item[Durability] The probability that stored data is not lost, commonly quoted as ``eleven nines''
    (99.999999999\%) for cloud object storage.
  \item[Expectation] A DLT data-quality rule with warn (\texttt{expect}), drop
    (\texttt{expect\_or\_drop}), or fail (\texttt{expect\_or\_fail}) semantics.
  \item[Feature Store] A governed repository of feature tables providing point-in-time-correct
    training sets and online serving; on Databricks, Feature Engineering in Unity Catalog.
  \item[Genie] Databricks AI/BI's natural-language analytics interface, generating SQL over curated
    datasets guided by instructions and example queries.
  \item[Lakehouse] An architecture that stores open table formats (such as Delta) in a data lake while
    offering warehouse-like SQL and governance; Databricks' founding architecture, also realized as
    Microsoft Fabric OneLake and AWS Lake Formation-governed lakes.
  \item[Liquid clustering] A Delta Lake layout technique (\texttt{CLUSTER BY}) that replaces fixed
    partition columns and Z-ORDER with clustering keys that can change without rewriting the table.
  \item[Materialized view] A precomputed, automatically maintained query result stored for fast reads,
    trading maintenance/storage cost for lower query latency (Databricks SQL, Snowflake, Redshift,
    BigQuery).
  \item[Medallion architecture] A progressive-refinement data layout with Bronze (raw), Silver
    (validated), and Gold (curated/enriched) layers.
  \item[Metastore] The Unity Catalog top-level container binding catalogs, credentials, and audit for
    a cloud region; the successor to the per-workspace Hive metastore.
  \item[Micro-partition] Snowflake's internal storage unit: an immutable, compressed, columnar file of
    roughly 50--500~MB of uncompressed data, annotated with per-column metadata for pruning.
  \item[MLflow] The open-source ML lifecycle platform (tracking, models, registry, evaluation);
    Databricks manages it and integrates the registry with Unity Catalog.
  \item[Mosaic AI] Databricks' AI suite: Foundation Model APIs, Vector Search, model serving, and
    evaluation tooling for GenAI and classic ML.
  \item[OPTIMIZE] The Delta Lake command that compacts small files, optionally applying
    \texttt{ZORDER BY} multi-dimensional clustering.
  \item[Partition pruning] Skipping storage units (micro-partitions, files, or partitions) that cannot
    match a query's predicates, reducing scanned data.
  \item[Photon] Databricks' vectorized C++ execution engine accelerating SQL and DataFrame workloads
    on clusters and SQL warehouses.
  \item[Point-in-time recovery (PITR)] Restoring a database or table to any moment within a retention
    window; realized as Delta Lake time travel/\texttt{RESTORE}, Snowflake Time Travel, and RDS PITR.
  \item[Serverless SQL warehouse] Databricks SQL compute that starts instantly, autoscales, and bills
    per second, removing cluster management for BI workloads.
  \item[Service principal] A non-human identity for automated jobs, pipelines, and CI/CD, granted
    least-privilege access; the Databricks counterpart to an AWS IAM role or Azure managed identity.
  \item[Small-file problem] The degradation caused by thousands of tiny data files: slow listings,
    metadata overhead, and poor scan efficiency; remediated by compaction (\texttt{OPTIMIZE}).
  \item[SQL warehouse] Databricks' SQL-optimized compute (serverless, pro, or classic) powering
    Databricks SQL queries, AI/BI dashboards, and Genie.
  \item[Streaming table] A DLT dataset backed by a Structured Streaming query, ingesting each source
    record exactly once; contrast with a DLT materialized view.
  \item[Structured Streaming] Spark's streaming engine, treating a stream as an unbounded table
    processed in micro-batches with checkpointed exactly-once semantics.
  \item[Time travel] Querying a Delta table as of a prior version or timestamp
    (\texttt{VERSION AS OF} / \texttt{TIMESTAMP AS OF}); Snowflake's analog is Time Travel.
  \item[Unity Catalog] Databricks' unified governance layer: metastore, catalogs, schemas, tables,
    volumes, functions, and models, with lineage, audit, and sharing.
  \item[VACUUM] The Delta Lake command that deletes data files no longer referenced by the transaction
    log's retention horizon (default 7 days), permanently ending time travel past that point.
  \item[Vector Search] Mosaic AI's managed vector index over Delta tables, powering semantic retrieval
    for retrieval-augmented generation (RAG) applications.
  \item[Volume] A Unity Catalog object governing access to non-tabular files in cloud storage.
  \item[Watermark] A Structured Streaming event-time threshold bounding how late events are accepted
    into windows and when state is evicted.
  \item[WiredTiger] MongoDB's default storage engine, using a configurable internal cache alongside the
    operating-system page cache, with journaling and checkpointing.
  \item[Workflow] A Databricks job composed of a DAG of tasks (notebooks, pipelines, SQL, dbt) with
    scheduling, retries, repair runs, and alerting.
  \item[Z-ordering] A Delta Lake optimization that co-locates related information in the same files by
    multi-dimensional clustering, improving data skipping; applied via \texttt{OPTIMIZE ... ZORDER BY}
    and largely superseded by liquid clustering for new tables.
\end{description}
